% ============================================================================
% LANGENTWURF LE-006: Arbeitsumgebung Mac
% Profil: S (Senioren 65+) | Tier: 2 (Standard)
% ============================================================================

\documentclass[11pt,a4paper]{article}
\usepackage[utf8]{inputenc}
\usepackage[T1]{fontenc}
\usepackage[ngerman]{babel}
\usepackage{geometry}
\usepackage{fancyhdr}
\usepackage{lastpage}
\usepackage{xcolor}
\usepackage{tcolorbox}
\usepackage{enumitem}
\usepackage{longtable}
\usepackage{hyperref}
\usepackage{amssymb}

\geometry{left=2.5cm, right=2.5cm, top=2.5cm, bottom=2.5cm}
\definecolor{fach}{HTML}{b35c00}
\definecolor{gray}{HTML}{666666}
\definecolor{successgreen}{HTML}{2a7a4b}

\tcbuselibrary{skins,breakable}
\newtcolorbox{scorebox}{ colback=white, colframe=fach, fonttitle=\bfseries, title=Didaktik-Score, breakable }

\pagestyle{fancy}
\fancyhf{}
\fancyhead[L]{\small\textcolor{gray}{Digitale Grundlagen -- 006 Arbeitsumgebung Mac}}
\fancyhead[R]{\small\textcolor{gray}{LE Tier 2 / Profil S}}
\fancyfoot[C]{\small\thepage{} / \pageref{LastPage}}
\renewcommand{\headrulewidth}{0.4pt}

\begin{document}

\begin{titlepage}
    \centering
    \vspace*{2cm}
    {\Large\textcolor{gray}{Digitale Grundlagen -- Senioren 65+}}\\[2cm]
    {\Huge\textcolor{fach}{\textbf{Langentwurf}}}\\[0.5cm]
    {\large Tier 2 | Profil S}\\[2cm]
    {\LARGE\textbf{006 -- Arbeitsumgebung Mac}}\\[0.5cm]
    {\large Modul A: Geräte \& Zugänge}\\[2cm]
    \begin{tabular}{rl}
        \textbf{Zeitrahmen:} & 75--90 Minuten \\
        \textbf{Vorkenntnisse:} & Grundverständnis Computer, Apple-ID \\
    \end{tabular}
    \vfill
    {\normalsize Stand: \today}
\end{titlepage}

\tableofcontents
\newpage

\section{Bedingungsanalyse}

\subsection{Ausgangssituation}
\begin{itemize}
    \item Mac-Bedienung unterscheidet sich deutlich von iPhone/iPad
    \item Maus/Trackpad statt Touchscreen
    \item Fenster-Konzept vs. Vollbild-Apps
    \item Teilnehmer kennen evtl. Windows -- Transfer möglich, aber auch Verwirrung
\end{itemize}

\subsection{Typische Probleme}
\begin{itemize}
    \item ``Wo ist der Start-Knopf?'' (Windows-Gewohnheit)
    \item ``Wie schließe ich ein Programm?''
    \item ``Wo sind meine Dateien?''
    \item ``Das Trackpad macht komische Sachen''
\end{itemize}

\section{Sachanalyse}

\subsection{Kernkonzepte}
\begin{enumerate}
    \item \textbf{Schreibtisch (Desktop):} Die Arbeitsfläche, wie ein echter Schreibtisch
    \item \textbf{Menüleiste:} Immer oben, wechselt je nach aktivem Programm
    \item \textbf{Dock:} Die Leiste unten mit den wichtigsten Programmen
    \item \textbf{Finder:} Der ``Aktenschrank'' -- Dateien und Ordner
    \item \textbf{Fenster:} Programme laufen in Fenstern (können verschoben werden)
    \item \textbf{Trackpad/Maus:} Klicken, Rechtsklick, Scrollen
\end{enumerate}

\subsection{Mac vs. Windows}
\begin{tabular}{|l|l|l|}
\hline
\textbf{Konzept} & \textbf{Windows} & \textbf{Mac} \\
\hline
Startmenü & Unten links & Gibt es nicht (Dock + Launchpad) \\
Dateiverwaltung & Explorer & Finder \\
Einstellungen & Systemsteuerung & Systemeinstellungen \\
Programme schließen & Rotes X & Rotes X schließt nur Fenster! \\
Rechtsklick & Rechte Maustaste & Zwei Finger tippen / Ctrl+Klick \\
\hline
\end{tabular}

\section{Didaktische Analyse}

\subsection{Stellung in der Reihe}
Letzte Einheit in Modul A. Optionale Vertiefung für Mac-Besitzer. Deutlich anders als iPhone/iPad.

\subsection{Besondere Herausforderung}
\begin{itemize}
    \item Mac-Bedienung ist \textbf{nicht} wie iPhone -- Transfer begrenzt
    \item Fenster-Konzept muss verstanden werden
    \item Trackpad-Gesten sind anders als Touchscreen-Gesten
\end{itemize}

\section{Lernziele}

\begin{tabular}{|c|p{7cm}|c|c|}
\hline
\textbf{Nr.} & \textbf{Die Teilnehmer können...} & \textbf{Stufe} & \textbf{Niveau} \\
\hline
FZ1 & Schreibtisch, Menüleiste und Dock benennen & Wissen & Basis \\
FZ2 & ein Programm aus dem Dock starten & Anwenden & Basis \\
FZ3 & ein Fenster verschieben und schließen & Anwenden & Basis \\
FZ4 & den Finder öffnen und Dateien finden & Anwenden & Basis \\
FZ5 & den Unterschied zwischen Fenster schließen und Programm beenden erklären & Verstehen & Vertiefung \\
\hline
\end{tabular}

\section{Verlaufsplanung}

\begin{longtable}{|p{1cm}|p{2cm}|p{5cm}|p{3cm}|p{2cm}|}
\hline
\textbf{Zeit} & \textbf{Phase} & \textbf{Inhalt} & \textbf{TN-Aktivität} & \textbf{Material} \\
\hline
0--10 & Einstieg & ``Was sehen Sie auf dem Bildschirm?'' Elemente benennen & Beobachten & -- \\
\hline
10--25 & Schreibtisch & Menüleiste (Apfel-Menü!), Dock erklären, Programm starten & Mitmachen & S.1 \\
\hline
25--40 & Finder & Der ``Aktenschrank'', Ordner öffnen, Dateien finden & Selbst navigieren & S.2 \\
\hline
40--55 & Fenster & Verschieben, Größe ändern, minimieren, schließen vs. beenden & Ausprobieren & S.2 \\
\hline
55--70 & Trackpad & Klicken, Rechtsklick, Scrollen, evtl. Gesten & Üben & S.2 \\
\hline
70--85 & Übungen & Checkliste durcharbeiten, typische Aufgaben & Bearbeiten & S.3 \\
\hline
\end{longtable}

\section{Lösungen}

\textbf{Wichtige Orte:}
\begin{itemize}
    \item Systemeinstellungen: Apfel-Menü $\rightarrow$ Systemeinstellungen
    \item Über diesen Mac: Apfel-Menü $\rightarrow$ Über diesen Mac
    \item Programm beenden: Apfel-Menü $\rightarrow$ [Programmname] beenden (oder Cmd+Q)
\end{itemize}

\section{Didaktik-Score}

\begin{scorebox}
\begin{tabular}{|c|l|c|c|p{3.5cm}|}
\hline
\# & Dimension & Gew. & Score & Notiz \\
\hline
1 & Differenzierung & 15\% & 7/10 & Komplexes Thema, wenig Stufen \\
2 & Taxonomie & 10\% & 9/10 & Anwenden dominiert \\
3 & Alltagsrelevanz & 10\% & 8/10 & Nur für Mac-Besitzer \\
4 & Aktivierung & 10\% & 9/10 & Viel Ausprobieren \\
5 & Scaffolding & 10\% & 8/10 & Schrittweise, aber komplex \\
6 & Feedback & 10\% & 8/10 & Checkliste vorhanden \\
7 & Sprachliche Klarheit & 10\% & 9/10 & Analogien helfen \\
8 & Motivation & 10\% & 8/10 & Erfolg bei Dateien finden \\
9 & Barrierefreiheit & 10\% & 8/10 & Trackpad kann schwierig sein \\
10 & Praktikabilität & 5\% & 7/10 & 85 Min ist lang \\
\hline
\multicolumn{3}{|r|}{\textbf{GESAMT:}} & \textbf{8.25/10} & Überarbeitung empfohlen \\
\hline
\end{tabular}
\end{scorebox}

\subsection{Verbesserungsvorschläge}
\begin{enumerate}
    \item \textbf{Praktikabilität:} In zwei Sitzungen aufteilen (Grundlagen + Finder)
    \item \textbf{Barrierefreiheit:} Maus als Alternative zum Trackpad anbieten
\end{enumerate}

\end{document}
